% AA-44 | Twierdzenie Liouville'a, pełna konieczność i dostateczność.
% Podstawa: AA-43, sources/pl/galois(2).tex; AA-45 dopiero po audytach.
\chapter{Twierdzenie Liouville'a o całkowaniu elementarnym}
\chaptermark{Twierdzenie Liouville'a}
\label{aa:ch:liouville-elementary-integration}

Funkcja mająca funkcję pierwotną nie musi mieć jej w klasie funkcji
elementarnych. Aby uczynić to pytanie ścisłym, funkcje zastępujemy
elementami ciała różniczkowego, a wyrażenia elementarne opisujemy
skończoną wieżą rozszerzeń. Narzędzia konstrukcyjne są już dostępne
w rozdziale~\ref{aa:ch:differential-fields}: jednoznaczne
przedłużenie derywacji do rozszerzenia algebraicznego, pochodne
śladu i normy oraz kryteria braku nowych stałych. Teraz
udowodnimy, że cała wieża może zostać wyeliminowana z
\emph{kryterium istnienia} elementarnej funkcji pierwotnej.

\section{Sformułowanie problemu i rola stałych}

\begin{definicja}[Postać Liouville'a]
\label{aa:def:liouville-form}
Niech $(K,D)$ będzie ciałem różniczkowym charakterystyki $0$,
$C=C_K$, a $f\in K$. Mówimy, że $f$ ma
\emph{postać Liouville'a w $K$}, jeżeli istnieją
$n\ge0$, $v\in K$, $u_1,\ldots,u_n\in K^\times$ oraz
$c_1,\ldots,c_n\in C$ takie, że
\begin{equation}\label{aa:eq:liouville-form}
 f=Dv+\sum_{j=1}^{n}c_j\frac{Du_j}{u_j}.
\end{equation}
Dla $n=0$ suma jest pusta. Współczynniki $c_j$
są \emph{stałymi derywacji}, a nie dowolnymi elementami $K$.
\end{definicja}

Równość~\eqref{aa:eq:liouville-form} nie zawiera
symboli logarytmów, dlatego ma sens w dowolnym
ciele różniczkowym charakterystyki $0$.
Jeżeli wybierzemy lokalne gałęzie logarytmów,
formalnym kandydatem na funkcję pierwotną jest
$v+\sum_j c_j\log u_j$. Trzeba jednak sprawdzić,
czy dołączanie logarytmów nie wprowadza nowych stałych.

\begin{przyklad}[Dwa elementarne przedstawienia]
\label{aa:prz:liouville-rational}
W $K=\mathbb C(x)$ z $Dx=1$ mamy
\[
 \frac{x}{x^2+1}=\frac12\frac{D(x^2+1)}{x^2+1},
 \qquad
 \frac1{x^2+1}
 =\frac1{2i}\left(\frac{D(x-i)}{x-i}
                   -\frac{D(x+i)}{x+i}\right).
\]
Odpowiednimi funkcjami pierwotnymi są lokalne gałęzie
$\frac12\log(x^2+1)$ oraz
$\frac1{2i}\log((x-i)/(x+i))$.
W drugiej tożsamości sprawdzamy jawnie
$\frac1{2i}((x-i)^{-1}-(x+i)^{-1})=(x^2+1)^{-1}$.
W obu przypadkach stałe współczynniki należą do $\mathbb C$.
\end{przyklad}

\section{Formalna część wielomianowa i pochodne logarytmiczne}

\begin{definicja}[Część wielomianowa w nieskończoności]
\label{aa:def:liouville-polynomial-part}
Niech $F$ będzie ciałem i $t$ elementem przestępnym nad $F$.
Każdy $R\in F(t)$ ma jednoznaczne przedstawienie
$R=P+Q$, gdzie $P\in F[t]$, a $Q=0$ albo
$Q=A/B$ dla $A,B\in F[t]$, $B\ne0$ i $\deg A<\deg B$.
Definiujemy $\operatorname{pp}_t(R)=P$.
\end{definicja}

Istnienie wynika z dzielenia z resztą licznika przez mianownik.
Dla jednoznaczności zauważmy, że niezerowy wielomian nie może
być właściwą funkcją wymierną: z $P=A/B$ wynikałoby
$\deg(PB)\ge\deg B>\deg A$.

Równoważnie, rozwijając mianownik w potęgach $t^{-1}$,
otrzymujemy jednoznaczne formalne rozwinięcie
$R=\sum_{j\le m}r_jt^j$ z $r_j\in F$,
gdzie tylko skończenie wiele wykładników jest nieujemnych.
Wówczas $\operatorname{pp}_t(R)=\sum_{j\ge0}r_jt^j$.
Jest to tożsamość w ciele formalnych szeregów Laurenta
$F((t^{-1}))$, a nie twierdzenie o granicy funkcji
analitycznej w nieskończoności.

\begin{lemat}[Części wielomianowe pochodnych]
\label{aa:lem:liouville-polynomial-parts}
Niech $(F,D)$ będzie ciałem różniczkowym
charakterystyki $0$, a $t$ elementem przestępnym nad $F$.
Załóżmy jeden z przypadków: $(\mathrm P)$ $Dt=a\in F$
albo $(\mathrm E)$ $Dt=ht$ dla $h\in F$.
W obu przypadkach dla $R\in F(t)$ zachodzi
\[
 \operatorname{pp}_t(DR)=D(\operatorname{pp}_t(R)).
\]
Jeśli $u\in F(t)^\times$ i
$u=b\,p/q$, gdzie $b\in F^\times$, a $p,q\in F[t]$
są unormowane, to
\begin{equation}\label{aa:eq:liouville-log-pp}
 \operatorname{pp}_t\left(\frac{Du}{u}\right)
 =\begin{cases}
       Db/b,&(\mathrm P),\\
       Db/b+(\deg p-\deg q)h,&(\mathrm E).
   \end{cases}
\end{equation}
W szczególności pochodna logarytmiczna nie wnosi
dodatnich potęg $t$.
\end{lemat}
\begin{proof}
Rozkład $R=P+Q$ z poprzedniej definicji pozwala
rozważyć osobno wielomian $P$ i właściwą funkcję
wymierną $Q=A/B$. W obu przypadkach $DP\in F[t]$,
bo $D(t^j)=ja t^{j-1}$ w $(\mathrm P)$ oraz
$D(t^j)=jht^j$ w $(\mathrm E)$.
Jeżeli $\deg A<\deg B=d$, to w obu przypadkach
$\deg DA\le\deg A$ (dla wielomianu stałego możliwe
jest $DA=0$), a $\deg DB\le d$.
Stąd licznik $(DA)B-A(DB)$ ma stopień co najwyżej
$\deg A+d<2d$, więc $DQ$ jest właściwe.
Dla $Q=0$ wniosek jest natychmiastowy.
To dowodzi pierwszej równości bez jakichkolwiek
założeń o rozkładzie $B$ na czynniki liniowe.

W drugiej części wykorzystujemy równość
$Du/u=Db/b+Dp/p-Dq/q$. Jeśli $p$ jest
unormowany stopnia $d$, to w przypadku $(\mathrm P)$
współczynnik przy $t^d$ w $Dp$ jest zerowy,
więc $\deg Dp<d$. W przypadku $(\mathrm E)$
współczynnik przy $t^d$ w $Dp$ wynosi $dh$,
więc $(Dp)/p-dh$ jest właściwą funkcją wymierną.
Analogiczne zdania zachodzą dla $q$; gdy
$p=1$ lub $q=1$, otrzymujemy wkład zerowy.
Część wielomianowa jest addytywna, skąd
wynika~\eqref{aa:eq:liouville-log-pp}.
\end{proof}

\begin{uwaga}[Dlaczego nie porównujemy dowolnych potęg Laurenta?]
\label{aa:uw:liouville-negative-powers}
Gdy $Dt=ht$, pochodna logarytmiczna
$D(t-1)/(t-1)=ht/(t-1)$ ma formalne rozwinięcie
$h+ht^{-1}+ht^{-2}+\cdots$. Nie wolno zatem
twierdzić, że pochodne logarytmiczne nie mają
ujemnych potęg. Udowodniony lemat dotyczy wyłącznie
\emph{części wielomianowej}, co dokładnie wystarczy
w dowodzie twierdzenia Liouville'a.
\end{uwaga}

\section{Trzy kroki schodzenia po wieży}

Zajmiemy się następującą sytuacją. Element $f$
należy do ciała niższego piętra $F$, natomiast
jego przedstawienie Liouville'a jest na razie znane
tylko w jednym rozszerzeniu $L/F$. Pokażemy, że
można je przenieść do $F$ bez zmiany ciała stałych.

\begin{lemat}[Schodzenie przez rozszerzenie algebraiczne]
\label{aa:lem:liouville-algebraic-descent}
Niech $(F,D)$ będzie ciałem różniczkowym
charakterystyki $0$ o ciele stałych $C$, a
$L/F$ rozszerzeniem skończonym algebraicznym
z przedłużoną derywacją. Jeśli $f\in F$ oraz
\[
 f=DV+\sum_{j=1}^{n}c_j\frac{DU_j}{U_j},
 \quad V\in L,\quad U_j\in L^\times,\quad c_j\in C,
\]
to $f$ ma postać Liouville'a w $F$.
\end{lemat}
\begin{proof}
Niech $d=[L:F]$. Ponieważ $\operatorname{char}F=0$,
liczba $d$ jest odwracalna w $C$.
Stosujemy ślad do obu stron tożsamości.
Ponieważ $f\in F$, mamy
$\operatorname{Tr}_{L/F}(f)=df$.
Na mocy twierdzenia~\ref{aa:tw:differential-trace-norm}
otrzymujemy
\[
 f=D\left(\frac{\operatorname{Tr}_{L/F}(V)}d\right)
   +\sum_{j=1}^{n}\frac{c_j}{d}
       \frac{D N_{L/F}(U_j)}{N_{L/F}(U_j)}.
\]
Ślad należy do $F$, a normy należą do $F^\times$.
Nowe współczynniki $c_j/d$ należą do $C$.
Nie jest potrzebne założenie normalności $L/F$.
\end{proof}

\begin{lemat}[Schodzenie przez krok logarytmiczny]
\label{aa:lem:liouville-log-descent}
Niech $(F,D)$ ma charakterystykę $0$ i ciało
stałych $C$. Niech $L=F(t)$, gdzie $t$ jest
przestępny nad $F$, $C_L=C$ oraz
$Dt=Db/b$ dla pewnego $b\in F^\times$.
Jeśli $f\in F$ ma postać Liouville'a w $L$
ze współczynnikami w $C$, to ma ją w $F$.
\end{lemat}
\begin{proof}
Zapiszmy
$f=DV+\sum_j c_j DU_j/U_j$ z $V\in F(t)$
i $U_j\in F(t)^\times$. W lemacie
\ref{aa:lem:liouville-polynomial-parts}
wybieramy przypadek $(\mathrm P)$ z $a=Db/b$.
Każde $U_j$ zapisujemy w postaci
$U_j=b_j p_j/q_j$, gdzie
$b_j\in F^\times$, a $p_j,q_j$ są unormowane.
Ponieważ $f\in F$, zastosowanie
$\operatorname{pp}_t$ daje
\begin{equation}\label{aa:eq:liouville-log-descent}
 f=DP+\sum_j c_j\frac{Db_j}{b_j},
 \quad P=\operatorname{pp}_t(V)\in F[t].
\end{equation}
W szczególności $DP\in F$.

Jeśli $P$ jest stały jako wielomian w $t$,
to~\eqref{aa:eq:liouville-log-descent} jest
już wymaganym przedstawieniem. W przeciwnym
razie niech $m=\deg_tP\ge1$ i
$P=p_mt^m+\cdots+p_0$, $p_m\ne0$.
Współczynnik przy $t^m$ w $DP$ wynosi
$Dp_m$; musi zatem być zerowy, a więc
$p_m\in C^\times$.
Gdyby $m\ge2$, współczynnik przy $t^{m-1}$
byłby równy $Dp_{m-1}+mp_m Db/b=0$.
Element $w=-p_{m-1}/(mp_m)\in F$
spełniałby $Dw=Db/b=Dt$.
Stąd $D(t-w)=0$, więc $t-w\in C_L=C\subseteq F$,
sprzecznie z przestępnością $t$ nad $F$.
Pozostaje $m=1$, zatem $P=p_0+ct$, $c\in C$.
Podstawiając do~\eqref{aa:eq:liouville-log-descent},
otrzymujemy
\[
 f=Dp_0+c\,\frac{Db}{b}
              +\sum_j c_j\frac{Db_j}{b_j}.
\]
Każdy składnik jest już zdefiniowany nad $F$.
\end{proof}

\begin{lemat}[Schodzenie przez krok wykładniczy]
\label{aa:lem:liouville-exp-descent}
Niech $(F,D)$ ma charakterystykę $0$ i ciało
stałych $C$. Niech $L=F(t)$, gdzie $t$ jest
przestępny nad $F$, $C_L=C$, $t\ne0$ oraz
$Dt/t=Da$ dla pewnego $a\in F$.
Jeśli $f\in F$ ma postać Liouville'a w $L$
ze współczynnikami w $C$, to ma ją w $F$.
\end{lemat}
\begin{proof}
Piszemy $f=DV+\sum_jc_j DU_j/U_j$ i
$U_j=b_jp_j/q_j$ z $b_j\in F^\times$ oraz
unormowanymi wielomianami $p_j,q_j$.
Lemat~\ref{aa:lem:liouville-polynomial-parts},
tym razem w przypadku $(\mathrm E)$ z $h=Da$,
daje
\begin{equation}\label{aa:eq:liouville-exp-descent}
 f=DP+\sum_jc_j\frac{Db_j}{b_j}+A\,Da,
 \qquad
 A=\sum_jc_j(\deg p_j-\deg q_j)\in C,
\end{equation}
gdzie $P=\operatorname{pp}_t(V)=\sum_{k=0}^{m}p_kt^k$.
Po prawej poza $DP$ wszystkie składniki należą do $F$,
więc $DP\in F$. Dla każdego $k\ge1$
współczynnik przy $t^k$ w $DP$ wynosi
$Dp_k+k(Da)p_k$ i musi być zerowy.
Jeśli $p_k\ne0$, to
$D(p_kt^k)=(Dp_k+k(Da)p_k)t^k=0$.
Wobec $C_L=C$ mamy $p_kt^k\in C^\times$,
czyli $t^k\in F$, co przeczy przestępności $t$.
Zatem $p_k=0$ dla wszystkich $k\ge1$,
a $P=p_0\in F$. Ponieważ $A\in C$,
mamy $A\,Da=D(Aa)$. Z~\eqref{aa:eq:liouville-exp-descent}
wynika więc
\[
 f=D(p_0+Aa)+\sum_jc_j\frac{Db_j}{b_j},
\]
czyli przedstawienie nad $F$.
\end{proof}

\section{Twierdzenie Liouville'a}

\begin{twierdzenie}[Liouville'a o całkowaniu elementarnym]
\label{aa:tw:liouville-elementary-integration}
Niech $(K,D)$ będzie ciałem różniczkowym
charakterystyki $0$ o ciele stałych $C=C_K$
i niech $f\in K$.
Następujące warunki są równoważne:
\begin{enumerate}
\item istnieje rozszerzenie elementarne
$E/K$ bez nowych stałych w sensie
Definicji~\ref{aa:def:elementary-differential-extension}
oraz element $F\in E$ taki, że $DF=f$;
\item $f$ ma postać Liouville'a w $K$, to znaczy
istnieją $n\ge0$, $v\in K$, $u_1,\ldots,u_n\in K^\times$
i $c_1,\ldots,c_n\in C$ spełniające
\eqref{aa:eq:liouville-form}.
\end{enumerate}
W szczególności nie trzeba przechodzić do rozszerzenia
algebraicznego, aby zapisać elementy $v,u_j$.
Nie zakładamy algebraicznej domkniętości $C$:
istotnym warunkiem jest \emph{brak nowych stałych}.
\end{twierdzenie}
\begin{proof}
\emph{Konieczność.} Załóżmy (1) i wybierzmy wieżę
elementarną $K=K_0\subseteq K_1\subseteq\cdots
\subseteq K_r=E$. Ponieważ
$C_K\subseteq C_{K_j}\subseteq C_E=C_K$,
wszystkie piętra mają to samo ciało stałych $C$.
Na najwyższym piętrze mamy przedstawienie
$f=DF$ odpowiadające~\eqref{aa:eq:liouville-form}
z $n=0$.

Przypuśćmy indukcyjnie, że $f$ ma postać
Liouville'a w $K_j$ ze współczynnikami w $C$.
Jeżeli generator $\theta_j$ jest algebraiczny
nad $K_{j-1}$, stosujemy lemat
\ref{aa:lem:liouville-algebraic-descent}.
Jeżeli generator jest przestępny i
$D\theta_j=Du/u$, stosujemy lemat
\ref{aa:lem:liouville-log-descent}.
Jeżeli generator jest przestępny,
$\theta_j\ne0$ i $D\theta_j/\theta_j=Dv$,
stosujemy lemat~\ref{aa:lem:liouville-exp-descent}.
Jeśli generator algebraiczny spełnia również
któryś z pozostałych warunków, używamy kroku
algebraicznego: nie ma potrzeby traktować
tych przypadków oddzielnie.
W każdym wariancie otrzymujemy przedstawienie
z elementami $K_{j-1}$ i współczynnikami w $C$.
Indukcja malejąca po $j=r,r-1,\ldots,1$
prowadzi do przedstawienia w $K$, co daje (2).

\emph{Dostateczność.} Załóżmy (2), zatem
$f=Dv+\sum_{j=1}^n c_j Du_j/u_j$.
Zbudujemy skończoną wieżę
$K=E_0\subseteq E_1\subseteq\cdots\subseteq E_n$,
zachowując ciało stałych $C$.
Na etapie $j$ kładziemy $a_j=Du_j/u_j\in K
\subseteq E_{j-1}$. Jeśli istnieje
$w_j\in E_{j-1}$ takie, że $Dw_j=a_j$,
nie powiększamy ciała, kładąc $E_j=E_{j-1}$.
Jeżeli taka funkcja pierwotna nie istnieje,
dołączamy \emph{przestępny} element $w_j$
z derywacją $Dw_j=a_j$ i kładziemy
$E_j=E_{j-1}(w_j)$. Jest to krok
logarytmiczny, bo $u_j\in E_{j-1}^\times$.
Lemat~\ref{aa:lem:primitive-extension-constants}
stwierdza, że w tej drugiej sytuacji
$C_{E_j}=C_{E_{j-1}}=C$. W pierwszej jest
to oczywiste.

Po wszystkich krokach $E=E_n$ jest
rozszerzeniem elementarnym bez nowych stałych.
Element $F=v+\sum_{j=1}^n c_jw_j\in E$
spełnia
$DF=Dv+\sum_jc_jDw_j
=Dv+\sum_jc_j Du_j/u_j=f$.
Dla $n=0$ wystarcza $E=K$ i $F=v$.
Dowodzi to (1).
\end{proof}

\begin{uwaga}[Co dokładnie zostało udowodnione?]
\label{aa:uw:liouville-scope}
Twierdzenie dotyczy istnienia funkcji pierwotnej
w \emph{elementarnym rozszerzeniu bez nowych stałych}
ustalonego ciała różniczkowego. Przy interpretacji
jako zarodków funkcji meromorficznych nad
$\mathbb C$ warunek braku nowych stałych jest
naturalny; gałęzie funkcji algebraicznych i
logarytmów wybieramy lokalnie. Jeżeli zezwalamy
na dowolne powiększanie ciała stałych, przed
zastosowaniem twierdzenia należy osobno uzasadnić
odpowiednią redukcję. Nie jest ona ukrytym
elementem niniejszego dowodu.
\end{uwaga}

\section{Kontrola twierdzenia na konkretnych przykładach}

\begin{przyklad}[Norma usuwa algebraiczny pierwiastek]
\label{aa:prz:liouville-norm-sqrt}
Weźmy $K=\mathbb C(x)$, $\alpha^2=x^2+1$,
$L=K(\alpha)$. Z rozdziału
\ref{aa:ch:differential-fields} wiemy, że
$D\alpha=x/\alpha$, zatem
\[
 \frac{D\alpha}{\alpha}
 =\frac{x}{x^2+1}\in K.
\]
W $L$ prawa strona jest pojedynczą
pochodną logarytmiczną.
Ponieważ $N_{L/K}(\alpha)=-(x^2+1)$,
uśrednienie śladu daje w $K$
\[
 \frac{x}{x^2+1}
 =\frac12\frac{D(-(x^2+1))}{-(x^2+1)}
 =\frac12\frac{D(x^2+1)}{x^2+1}.
\]
Znak normy nie wpływa na pochodną logarytmiczną.
Przykład ilustruje rzeczywistą zmianę elementu
$u_j$ przy schodzeniu po kroku algebraicznym.
\end{przyklad}

\begin{przyklad}[Krok logarytmiczny]
\label{aa:prz:liouville-log-step}
Niech $K=\mathbb C(x)$, $Dx=1$, a $t$ będzie
przestępny z $Dt=1/x$.
Fakt, że $K(t)$ nie ma nowych stałych, wynika
z rozkładu funkcji wymiernych na ułamki proste
i lematu~\ref{aa:lem:primitive-extension-constants}.
Dla $V=t^2+x$ zachodzi $DV=2t/x+1$,
co nie należy do $K$; człon $t^2$ nie może
zostać ukryty w pochodnych logarytmicznych,
gdyż ich części wielomianowe są stałe w $t$.
Natomiast dla $V=3t+x^2$
mamy $DV=3/x+2x\in K$ i twierdzenie
sprowadza $DV$ do postaci
$D(x^2)+3\,Dx/x$.
W ten sposób obliczenie współczynników
przy potęgach $t$ wyjaśnia oba przypadki.
\end{przyklad}

\begin{przyklad}[Krok wykładniczy]
\label{aa:prz:liouville-exp-step}
Niech $K=\mathbb C(x)$, $Dx=1$ oraz
$t=e^x$ rozumiane jako przestępny generator
z $Dt=t$. Brak nowych stałych wykazano
w przykładzie~\ref{aa:prz:exponential-extension}.
Dla $u=xt^2/(t+1)$ otrzymujemy
\[
 \frac{Du}{u}=\frac1x+2-\frac{t}{t+1}.
\]
Ponieważ $t/(t+1)=1-1/(t+1)$,
część wielomianowa tego wyrażenia wynosi
$1/x+1$, zgodnie ze wzorem
\eqref{aa:eq:liouville-log-pp}:
$\deg(xt^2)-\deg(t+1)=1$ po wyłączeniu $x$.
W lemacie wykładniczym otrzymana
stała $1$ jest dokładną pochodną $Dx$.
\end{przyklad}

\begin{przyklad}[Dlaczego zabronione są nowe stałe?]
\label{aa:prz:liouville-new-constants}
Jeżeli do $K=\mathbb C(x)$ dołączymy
przestępny element $t$ z $Dt=1=Dx$,
to $D(t-x)=0$ mimo $t-x\notin K$.
Przykład ten nie stanowi kroku
logarytmicznego bez nowych stałych:
kryterium z AA-43 jest naruszone.
W szczególności argument, że
$D(p_mt^m)=0$ zmusza $p_mt^m$ do należenia
do pierwotnego ciała stałych, byłby wtedy błędny.
\end{przyklad}

\section{Granice wyniku i następny krok}

Twierdzenie Liouville'a sprowadza problem
całkowalności elementarnej do znalezienia
elementów $v,u_j\in K$ i stałych $c_j$.
Samo w sobie nie rozstrzyga jeszcze,
czy dane $f$ posiada takie przedstawienie.
W szczególności nie wolno przyjąć dowolnej
funkcji pierwotnej w kroku elementarnym,
ponieważ wtedy problem stałby się tautologiczny.
W kolejnym rozdziale zastosujemy uzyskane
kryterium do funkcji hiperwykładniczych:
porównując współczynniki przy dodatnich
potęgach przestępnego generatora,
sprowadzimy pytanie o elementarność konkretnych
całek do równań różniczkowych w $\mathbb C(x)$.
